% Where the single writer's cycles go: P1 sec:measured, "Where the writer's
% cost goes" -- cycles per update on the writer's own thread, one writer, ONE
% READER, nodes in walk order, no writer mutex (the ratio there is 1.5), engine
% 2793224e.  The numbers are P1's text (from the benchmark tree's
% scripts/p1_writer_profile_ret.csv, not shipped in P1's data/), so they are
% typed here.
%
% Two bars on ONE cycle scale, the shared categories first and in the same
% order, so what the transacted writer adds sticks out on the right: the commit
% and the descriptor slab are 19 + 95 = 114 cycles, which is also the whole
% difference (243 - 129), because the three shared categories net to zero
% (+23 write function, +6 call_rcu, -29 node allocation).  P1 does not claim
% the -29: it says node allocation "reads lower" and that the reason is not
% established.  So the slide shows the two bars and draws no brace over a
% "gap share".
%
% Sampled at retirement (plain cycles), not with the precise event: on this
% machine the precise event tags an op as it enters the pipeline and charges a
% stall to whatever runs next, which smeared the slab's cost over the write
% function in the earlier version of this figure.
%
% What each category holds follows the categorization script
% (efficios-trie-benchmark scripts/p1_writer_profile_categorize.py), by symbol
% self time: "commit" is urcu_txn_sw_commit_flavor, into which install, flip and
% settle are inlined; "write function" is the harness's write function and
% loop plus the engine's list calls, with the append now inlined into them.
%
% Until 2026-10-03 this figure was a single bar of "shares of the gap" at one
% writer and NO reader (46% call_rcu, 20% slab, 21% staging, 11% commit, engine
% 18809ea8).  That configuration is gone from P1: plain RCU's reclaim does not
% keep up there, so the gap divided by a rate it cannot sustain.
\begin{tikzpicture}[x=0.555mm, y=1cm,
  seg/.style = {draw=paper, line width=0.8pt},
  num/.style = {font=\small\bfseries, text=paper},
  nam/.style = {font=\footnotesize, text=ink, anchor=north, align=center},
  row/.style = {font=\footnotesize, text=ink, anchor=south west},
]
% ---- rculist: 129 cycles
\node[row] at (0,1.86) {\code{rculist}\enspace{\color{muted}129 cycles per update}};
\fill[seg, fill=stripes]  (0,1.24)  rectangle (47,1.84);
\fill[seg, fill=muted]    (47,1.24) rectangle (79,1.84);
\fill[seg, fill=onelock]  (79,1.24) rectangle (129,1.84);
\node[num, text=ink] at (23.5,1.54) {47};   % paper on this blue is < 4.5:1
\node[num] at (63,1.54)   {32};
\node[num] at (104,1.54)  {50};

% ---- txn_sw_list: 243 cycles
\node[row] at (0,0.82) {\code{txn\_sw\_list}\enspace{\color{muted}243 cycles per update}};
\fill[seg, fill=stripes]  (0,0.2)   rectangle (70,0.8);
\fill[seg, fill=muted]    (70,0.2)  rectangle (108,0.8);
\fill[seg, fill=onelock]  (108,0.2) rectangle (129,0.8);
\fill[seg, fill=accent]   (129,0.2) rectangle (148,0.8);
\fill[seg, fill=old]      (148,0.2) rectangle (243,0.8);
\node[num, text=ink] at (35,0.5)    {70};
\node[num] at (89,0.5)    {38};
\node[num] at (118.5,0.5) {21};
\node[num] at (138.5,0.5) {19};
\node[num] at (195.5,0.5) {95};

\node[nam] at (35,0.12)    {write function\\{\scriptsize\color{muted}list op + record}};
\node[nam] at (89,0.12)    {\code{call\_rcu}};
\node[nam] at (118.5,0.12) {nodes};
\node[nam] at (138.5,0.12) {\textbf{commit}};
\node[nam] at (195.5,0.12) {txn descriptor slab\\{\scriptsize\color{muted}allocate, hand back; four locked instructions}};
\end{tikzpicture}
